\documentclass[11pt]{article}
\usepackage[T1]{fontenc}
\usepackage{lmodern}
\usepackage[margin=1.05in]{geometry}
\usepackage{amsmath,amssymb,amsthm}
\usepackage{booktabs}
\usepackage{microtype}
% MiKTeX's ORCID icon font is not scalable; protrusion remains enabled while
% font expansion is disabled to keep pdfTeX builds portable and deterministic.
\microtypesetup{expansion=false}
\usepackage{xcolor}
\definecolor{linknavy}{RGB}{25,55,125}
\usepackage[colorlinks=true,linkcolor=linknavy,citecolor=linknavy,urlcolor=linknavy]{hyperref}
\hypersetup{
  pdftitle={Frobenius Minimality of the Huneke--Wiegand Counterexample in Symmetric Numerical Semigroup Rings},
  pdfauthor={Felipe Santibanez-Leal},
  pdfsubject={Certified Frobenius-minimality theorem for a two-generated monomial-ideal counterexample},
  pdfkeywords={Huneke-Wiegand conjecture, numerical semigroups, SAT, DRAT, computer-assisted proof}
}
\usepackage{orcidlink}
\usepackage{fancyhdr}

\newtheorem{theorem}{Theorem}[section]
\newtheorem{proposition}[theorem]{Proposition}
\newtheorem{lemma}[theorem]{Lemma}
\theoremstyle{definition}
\newtheorem{definition}[theorem]{Definition}
\theoremstyle{remark}
\newtheorem{remark}[theorem]{Remark}

\newcommand{\N}{\mathbb{N}}
\newcommand{\kfield}{\Bbbk}
\newcommand{\one}{\mathbf{1}}

\title{Frobenius Minimality of the Huneke--Wiegand Counterexample\\
in Symmetric Numerical Semigroup Rings}
\author{Felipe Santiba\~nez-Leal\,\orcidlink{0000-0002-0150-3246}\\
\small CAOS Research program, \texttt{github.com/fsantibanezleal/CAOS\_RESEARCH}}
\date{2026-08-02\\[2pt]\normalsize Version 0.01}

\pagestyle{fancy}
\fancyhf{}
\fancyhead[L]{\small CAOS Research Preprint}
\fancyhead[R]{\small Huneke--Wiegand Frobenius Minimality \textperiodcentered\ v0.01}
\fancyfoot[C]{\small \thepage}
\renewcommand{\headrulewidth}{0.4pt}
\fancypagestyle{plain}{\fancyhf{}\fancyfoot[C]{\small \thepage}\renewcommand{\headrulewidth}{0pt}}

\begin{document}
\maketitle
\begin{center}\small
\fbox{\parbox{0.92\textwidth}{\centering
\textbf{PREPRINT} \textperiodcentered\ not peer reviewed \textperiodcentered\ CAOS Research programme\\
First paper of the Huneke--Wiegand extension record \textperiodcentered\ version 0.01 \textperiodcentered\ 2026-08-02\\
Concept DOI \href{https://doi.org/10.5281/zenodo.21763582}{10.5281/zenodo.21763582} (resolves to latest) \textperiodcentered\ this version: \href{https://doi.org/10.5281/zenodo.21763583}{10.5281/zenodo.21763583}\\
Licensed CC BY 4.0 \textperiodcentered\ sources and machine records: \texttt{github.com/fsantibanezleal/CAOS\_RESEARCH}\\[2pt]
\textbf{Provenance.} Son Pham discovered the counterexample; Professor Craig Huneke independently
verified it. This paper contributes a separate certified Frobenius-minimality result.
}}
\end{center}

\begin{abstract}
Son Pham publicly identified a counterexample to the Huneke--Wiegand conjecture in the class of
two-generated monomial ideals over symmetric numerical semigroup rings, subsequently verified by
Professor Craig Huneke. We preserve that discovery priority and answer a different question: what
is the least Frobenius number at which such a counterexample can occur? We prove by proof-carrying
exact computation that the answer is $181$, attained by Pham's semigroup at ideal shift $14$.
First, an implementation of the complete Blanco--Rosales tree exhausts all $48{,}954$ symmetric
semigroups with Frobenius number below $69$ and all $1{,}503{,}391$ associated gap cases. An
independent Boolean route checks $1{,}156$ fixed-pair UNSAT certificates for the same range. We then
introduce a one-hot shift-selector CNF that asks one existence question per Frobenius number. Its
DRAT certificates exclude every odd value from $69$ through $179$, while the formula at $181$ is
satisfiable and decodes exactly to the public semigroup and shift. A separately generated fixed-pair
formula returns the same model. An independent audit rehashes all $228$ files in the minimality
search, totaling $760{,}081{,}739$ bytes, rechecks both SAT models, and reproduces the strict-order
manifest with no mismatch. Complete theorem-tree enumeration at $F=69,71,73,75$ provides an
object-level cross-check on another $79{,}790$ semigroups and $2{,}933{,}163$ gap cases. The theorem
is scoped to symmetric numerical semigroup rings and nonprincipal two-generated monomial ideals. It
does not assert uniqueness, minimum multiplicity or embedding dimension, or minimality for arbitrary
Gorenstein domains and modules.
\end{abstract}

\section{Introduction}

Let $R$ be a one-dimensional Gorenstein local domain and let $M$ be a nonzero finitely generated
nonfree module. The Huneke--Wiegand conjecture predicted nontrivial torsion in
$M\otimes_R\operatorname{Hom}_R(M,R)$ \cite{HunekeWiegand1994}. For a numerical semigroup
$\Gamma\subseteq\N$, the semigroup ring $\kfield[t^\Gamma]$ is Gorenstein exactly when $\Gamma$ is
symmetric. This specialization converts part of the module problem into finite additive
combinatorics. Garc\'ia-S\'anchez and Leamer proved broad positive classes and reported a
NumericalSgps check for every symmetric semigroup with Frobenius number below $69$
\cite{GarciaSanchezLeamer2013}.

In 2026, Son Pham published a symmetric numerical semigroup and a two-generated monomial ideal for
which the predicted torsion is absent \cite{Pham2026}. Professor Craig Huneke supplied an
independent verification archived with the public candidate record. That example has Frobenius
number $181$. The discovery settles the original conjecture negatively, but does not by itself say
whether $181$ is small, minimal, or typical.

This paper reports a computer-assisted minimality theorem. Claims are labeled \textbf{[MV]} when
they are verified by exact machine computation and independently checked artifacts, and
\textbf{[D]} when derived in the text. No conjectural claim is needed for the main theorem.

\paragraph{Contributions.}
\begin{enumerate}
  \item We give a self-contained finite-window reduction and a direct proof that the selector CNF is
  equisatisfiable with existence of a rigid two-generated monomial ideal.
  \item We reproduce the published $F<69$ boundary by two independent complete routes: the
  Blanco--Rosales tree and checked DRAT proofs.
  \item We certify UNSAT for every odd $F=69,71,\ldots,179$, then recover the exact public model at
  $(F,s)=(181,14)$.
  \item We independently rehash and inspect the complete search artifact set, and cross-check the
  first four new odd Frobenius values by explicit theorem-tree enumeration.
  \item We identify two structural boundaries: the public semigroup lies outside the generalized
  arithmetic-sequence positive family, and the visible fixed-offset generator blocks do not extend
  to a naive parametric family.
\end{enumerate}

\section{The finite semigroup problem}

\begin{definition}
A \emph{numerical semigroup} is a cofinite additive submonoid $\Gamma\subseteq\N$. Its Frobenius
number is $F(\Gamma)=\max(\N\setminus\Gamma)$. It is \emph{symmetric} when
\[
  n\in\Gamma \quad\Longleftrightarrow\quad F(\Gamma)-n\notin\Gamma
  \qquad (0\le n\le F(\Gamma)).
\]
\end{definition}

Fix a symmetric semigroup $\Gamma$ with Frobenius number $F$ and a positive gap
$s\notin\Gamma$. Normalize the corresponding two-generated relative ideal to $(0,s)$. Define
\begin{align*}
E_s &= \{n\in\N:n,n+s\in\Gamma\},\\
D_s &= \{n\in\N:n,n+s,n+2s\in\Gamma\}.
\end{align*}
The Garc\'ia-S\'anchez--Leamer reduction identifies the counterexample condition with
\begin{equation}\label{eq:rigidity}
  D_s=E_s+E_s.
\end{equation}
Equivalently, every length-two arithmetic sequence of step $s$ in $\Gamma$ decomposes as a sum of
two length-one sequences \cite{GarciaSanchezLeamer2013,Landeros2024}.

\begin{lemma}[Automatic inclusion, {[D]}]\label{lem:auto}
For every numerical semigroup $\Gamma$ and every positive integer $s$,
$E_s+E_s\subseteq D_s$.
\end{lemma}
\begin{proof}
If $a,b\in E_s$, then $a,a+s,b,b+s\in\Gamma$. Closure gives
$a+b\in\Gamma$, $a+b+s=a+(b+s)\in\Gamma$, and
$a+b+2s=(a+s)+(b+s)\in\Gamma$. Hence $a+b\in D_s$.
\end{proof}

\begin{lemma}[Finite window, {[D]}]\label{lem:window}
Equation~\eqref{eq:rigidity} holds if and only if
$D_s\cap[0,2F+1]\subseteq E_s+E_s$.
\end{lemma}
\begin{proof}
Only the reverse inclusion needs checking by Lemma~\ref{lem:auto}. Every integer at least $F+1$
belongs to $\Gamma$, hence to $E_s$ and $D_s$. Let $m=\min E_s$; then $m\le F+1$. For every
$n\ge m+F+1$, both $m$ and $n-m$ lie in $E_s$, so $n\in E_s+E_s$. Since
$m+F+1\le2F+2$, every $n\ge2F+2$ is automatic. The displayed finite window is therefore complete.
\end{proof}

\begin{lemma}[Parity, {[D]}]\label{lem:parity}
The Frobenius number of a symmetric numerical semigroup is odd.
\end{lemma}
\begin{proof}
If $F$ were even, symmetry at $n=F/2$ would require $F/2$ to belong to $\Gamma$ exactly when it
does not belong to $\Gamma$.
\end{proof}

\section{The public model and independent calibration}

Pham's public semigroup is
\begin{align*}
\Gamma=\langle&56,57,58,63,64,70,71,72,73,74,75,76,77,78,79,80,81,82,83,\\
&87,89,90,93,95,96,97\rangle,
\end{align*}
with the ideal generated by $t^{56}$ and $t^{70}$, hence normalized shift $s=14$
\cite{Pham2026}. An independent Singular/4ti2 calculation and a separate standard-library finite
route verify $F=181$, conductor $182$, genus $91$, symmetry, and equality~\eqref{eq:rigidity}
\textbf{[MV]}. The Singular route uses a $322$-generator toric standard basis; both reduced ideal
differences vanish. The hypersurface control $\langle4,5\rangle$ rejects equality with explicit
residues \textbf{[MV]}.

The endomorphism overring of the ideal has value semigroup
\[
  v(\operatorname{End}_R(I))=\Gamma\cup\{101,107,181\},
\]
with Frobenius number $125$, genus $88$, and type $24$ \textbf{[MV]}. This explains why positive
criteria that require a Gorenstein endomorphism ring do not apply, but it does not imply minimality.

\section{Complete reproduction below 69}

For fixed odd $F$, irreducible numerical semigroups are exactly the symmetric ones. Blanco and
Rosales construct all of them as a rooted tree \cite{BlancoRosales2013}. Starting from
\[
 C(F)=\{0,(F+1)/2,(F+1)/2+1,\ldots\}\setminus\{F\},
\]
a child replaces a suitable minimal generator $x>F/2$ by $F-x$. Their theorem supplies exact
child conditions, completeness, and a unique parent.

Our direct implementation validates every node for closure, symmetry, genus, minimal generators,
and its theorem parent. It then checks every gap by Lemma~\ref{lem:window} and retains a first
element of $D_s\setminus(E_s+E_s)$ whenever equality fails. The full range $F=1,3,\ldots,67$
contains $48{,}954$ semigroups and $1{,}503{,}391$ gap cases; none is rigid \textbf{[MV]}. At
$F=11$ the implementation returns the six nodes listed by Blanco--Rosales, while an invalid
mutation is rejected by closure at $(1,1)$.

An independent route generates one fixed-pair DIMACS formula for each $(F,s)$, including shifts
that are forced not to be gaps. CaDiCaL emits a DRAT proof for every UNSAT result, and DRAT-trim
checks each proof against the exact CNF \cite{Cadical2023,DRAT2014}. All $1{,}156$ formulas are
UNSAT and independently verified \textbf{[MV]}. A committed auditor rehashes $4{,}624$ external
files totaling $211{,}671{,}241$ bytes with zero mismatch.

\section{One selector formula per Frobenius value}

The fixed-pair route duplicates the same semigroup constraints once per shift. We instead introduce
Boolean membership variables $h_n$ for $0\le n\le F$ and one-hot variables $q_s$ for
$1\le s\le F$.

\subsection{Semigroup and selector clauses}

The CNF contains units $h_0$ and $\neg h_F$, the symmetry clauses
\[
 (h_n\lor h_{F-n})\land(\neg h_n\lor\neg h_{F-n}),
\]
and closure clauses $\neg h_a\lor\neg h_b\lor h_{a+b}$ for $a+b\le F$. The selectors satisfy one
at-least-one clause, all pairwise at-most-one clauses, and
$q_s\Rightarrow\neg h_s$. Thus every model chooses exactly one gap.

For $0\le n\le2F+1$, guarded Tseitin clauses impose, under the selected $q_s$,
\begin{align*}
 e_n &\longleftrightarrow h_n\land h_{n+s},\\
 d_n &\longleftrightarrow h_n\land h_{n+s}\land h_{n+2s},
\end{align*}
where $h_j$ is the constant true for $j>F$. For each unordered decomposition $n=a+b$, a variable
$p_{n,a,b}$ is equivalent to $e_a\land e_b$. Finally,
\[
 d_n\Rightarrow\bigvee_{a+b=n}p_{n,a,b}.
\]
The automatic inclusion from Lemma~\ref{lem:auto} need not be duplicated in the formula.

\begin{proposition}[Selector equisatisfiability, {[D]}]\label{prop:cnf}
For positive odd $F$, the selector CNF is satisfiable if and only if there exist a symmetric
numerical semigroup $\Gamma$ with Frobenius number $F$ and a gap $s$ satisfying
$D_s=E_s+E_s$.
\end{proposition}
\begin{proof}
The semigroup clauses decode exactly a symmetric semigroup of Frobenius number $F$, and the one-hot
clauses decode a unique gap. Because exactly one guard is true, all $e_n$ and $d_n$ equal the
definitions for that gap. The $p$ equivalences and final disjunctions impose
$D_s\subseteq E_s+E_s$ through $2F+1$. Lemmas~\ref{lem:auto} and~\ref{lem:window} give global
equality. Conversely, from a rigid pair assign $h$ and $q$ from the pair, assign $e,d$ by their
definitions, and assign each $p$ by its conjunction. Every clause is then true.
\end{proof}

The selector encoding passes two adversarial calibrations \textbf{[MV]}. Every odd $F\le67$ is
UNSAT with an accepted proof, while $F=181$ is SAT. The latter formula has $34{,}397$ variables and
$363{,}912$ clauses, selects $s=14$, and decodes exactly to the public membership vector. A newly
generated fixed-pair formula at $(181,14)$ returns the same vector.

\section{The certified search and minimality theorem}

We scan odd $F$ in strict increasing order. Table~\ref{tab:search} summarizes the closed run.

\begin{table}[ht]
\centering
\caption{Selector-CNF minimality search.}
\label{tab:search}
\begin{tabular}{lr}
\toprule
Quantity & Value\\
\midrule
Odd Frobenius values requested & 57 ($69$ through $181$)\\
UNSAT with accepted DRAT proof & 56 ($69$ through $179$)\\
Validated SAT models & 1 ($181$, shift $14$)\\
UNKNOWN or timeout & 0\\
Wall time & $4{,}303.86$ seconds\\
CNF bytes & $169{,}745{,}433$\\
DRAT proof bytes & $586{,}785{,}257$\\
Slowest solve & $F=175$, $218.83$ seconds\\
\bottomrule
\end{tabular}
\end{table}

\begin{theorem}[Frobenius minimality, {[MV+D]}]\label{thm:minimal}
Among symmetric numerical semigroups admitting a nonprincipal two-generated monomial ideal that is
a counterexample to the Huneke--Wiegand prediction, the least Frobenius number is $181$. It is
attained by Pham's public example at shift $14$.
\end{theorem}
\begin{proof}
By Lemma~\ref{lem:parity}, only odd Frobenius values occur. Complete tree enumeration and the
independent fixed-pair proof suite exclude every odd $F\le67$. Proposition~\ref{prop:cnf} and the
accepted selector DRAT proofs exclude every odd $F$ from $69$ through $179$. The validated
selector and fixed-pair models at $(181,14)$ provide existence. These cases exhaust all smaller
positive Frobenius values.
\end{proof}

The search audit independently rehashes all $228$ expected CNF, proof, solver-log, and checker-log
files, totaling $760{,}081{,}739$ bytes. It also reconstructs both SAT membership masks, validates
closure, symmetry, and rigidity, checks strict odd ordering, and recomputes aggregate SHA-256
\[
\texttt{0f580de2707a00fdd52e1b3c04e7767b97ce7b0a826593b119e9a49ae04da743}.
\]
No file is missing, no hash differs, and no status marker or semantic check fails \textbf{[MV]}.

\section{Independent object-level checks beyond the old frontier}

The theorem-tree route was continued beyond $67$ without consulting SAT models. Table~\ref{tab:tree}
lists the first four new odd values.

\begin{table}[ht]
\centering
\caption{Complete Blanco--Rosales tree cross-checks beyond the published boundary.}
\label{tab:tree}
\begin{tabular}{rrrrr}
\toprule
$F$ & Symmetric semigroups & Gap cases & Rigid cases & Seconds\\
\midrule
69 & 11,276 & 394,660 & 0 & 137.00\\
71 & 19,812 & 713,232 & 0 & 235.37\\
73 & 25,405 & 939,985 & 0 & 333.21\\
75 & 23,297 & 885,286 & 0 & 330.10\\
\midrule
Total & 79,790 & 2,933,163 & 0 & 1,035.68\\
\bottomrule
\end{tabular}
\end{table}

This is not needed to trust a DRAT derivation, but it attacks the custom encoding at the object
level and supplies explicit failed-rigidity witnesses for every tested gap. The routes agree.

\section{Structural boundaries and a failed family ansatz}

Landeros et al. prove the positive result for semigroups generated by a generalized arithmetic
sequence $\langle a,ah+d,\ldots,ah+kd\rangle$ \cite{Landeros2024}. The public semigroup cannot lie
in this class \textbf{[D]}. Its multiplicity forces $a=56$. Since $57\in\Gamma$ and no sum of two
positive semigroup elements is below $112$, the presentation forces $h=d=1$. Since
$63\in\Gamma$, it would then include $59,60,61,62$, but all four are gaps.

The seed also displays $56=4s$ and $181=13s-1$ at $s=14$, with generator blocks
\[
4s+\{0,1,2,7,8\},\qquad 5s+[0,s-1],\qquad
6s+\{3,5,6,9,11,12,13\}.
\]
A predeclared exact sweep refutes the naive fixed-offset lift \textbf{[MV]}. Only $s=14$ passes
among even $s\le100$. At $s=16$, the generated semigroup has $F=205$ rather than $207$; at
$s=28$ the predicted Frobenius value returns but symmetry fails. This negative result prevents the
minimality theorem from being overinterpreted as evidence of an immediate family.

\section{Reproducibility and trust boundary}

All compact code, hypotheses, manifests, audits, and verdicts are in the public repository under
\begin{center}
\url{problems/commutative-algebra/huneke-wiegand/}.
\end{center}
Heavy CNFs and proofs are retained outside Git with hashes committed in the manifests. The proof
toolchain is pinned as follows:
\begin{center}
\begin{tabular}{@{}lp{0.65\textwidth}@{}}
\toprule
Component & Identity\\
\midrule
CaDiCaL & reported 1.7.3; Ubuntu package 1.7.4-1\\
CaDiCaL binary SHA-256 & {\fontsize{7}{8}\selectfont\ttfamily 7b73df0a6d9cf3c751a1948300e5baff8e82c4d39bcd88f0c063b5f5cfb8b33e}\\
DRAT-trim commit & {\fontsize{7}{8}\selectfont\ttfamily 2e3b2dc0ecf938addbd779d42877b6ed69d9a985}\\
DRAT-trim binary SHA-256 & {\fontsize{7}{8}\selectfont\ttfamily 92f0aa9575ed519d66a99b8b1b3dde6ece4618ae4c202a3a4b200265dda0aa7a}\\
\bottomrule
\end{tabular}
\end{center}

The trusted base comprises the cited tree theorem, the finite reduction proved here, the small
standard-library model checker, and DRAT-trim's verification of each proof against its exact CNF.
The SAT solver need not be trusted for UNSAT. For SAT, the independently extracted model is checked
without consulting solver internals.

\section{Scope, attribution, and open questions}

Son Pham retains discovery priority for the counterexample. Professor Craig Huneke's verification
is independent external evidence. The contribution proved here is Frobenius minimality in the
specified semigroup/ideal class and its reproducible certificate suite.

Theorem~\ref{thm:minimal} does not prove uniqueness at $F=181$, minimum multiplicity, minimum
embedding dimension, minimum shift, or minimality among arbitrary Gorenstein domains and modules.
The following questions remain open:
\begin{enumerate}
  \item How many nonisomorphic rigid pairs occur at $F=181$?
  \item Is there a stronger minimum under multiplicity, embedding dimension, or number of minimal
  generators?
  \item Does a nontrivial infinite family exist after the fixed-offset ansatz fails?
  \item Which additional hypotheses on the endomorphism ring recover a positive theorem while
  excluding the public example?
  \item Can the finite certificates be imported into a smaller formally verified checker?
\end{enumerate}

\section*{Acknowledgments}

The author thanks Son Pham for making the candidate and verification materials public, and Craig
Huneke for the independent verification recorded with that release. Neither is an author of this
paper, and this acknowledgment does not imply endorsement of the minimality result.

{\small
\begin{thebibliography}{99}

\bibitem{HunekeWiegand1994}
C. Huneke and R. Wiegand,
\emph{Tensor products of modules and the rigidity of Tor},
Math. Ann. 299 (1994), 449--476.
\href{https://doi.org/10.1007/BF01459793}{doi:10.1007/BF01459793}.

\bibitem{GarciaSanchezLeamer2013}
P. A. Garc\'ia-S\'anchez and M. J. Leamer,
\emph{Huneke--Wiegand conjecture for complete intersection numerical semigroup rings},
J. Algebra 391 (2013), 114--124.
\href{https://doi.org/10.1016/j.jalgebra.2013.06.007}{doi:10.1016/j.jalgebra.2013.06.007}.

\bibitem{BlancoRosales2013}
V. Blanco and J. C. Rosales,
\emph{The tree of irreducible numerical semigroups with fixed Frobenius number},
Forum Math. 25 (2013), 1249--1261.
\href{https://doi.org/10.1515/form.2011.151}{doi:10.1515/form.2011.151}.

\bibitem{Landeros2024}
M. Landeros, C. O'Neill, R. Pelayo, K. Pe\~na, J. Ren, and B. Wissman,
\emph{Families of numerical semigroups and a special case of the Huneke--Wiegand conjecture},
Commun. Algebra 52 (2024), 5196--5206.
\href{https://doi.org/10.1080/00927872.2024.2362934}{doi:10.1080/00927872.2024.2362934}.

\bibitem{Pham2026}
S. Pham,
\emph{Huneke--Wiegand candidate verification}, public software and certificate repository (2026).
\url{https://github.com/sonpham-org/huneke-wiegand-candidate-verification}.

\bibitem{Cadical2023}
A. Biere, K. Fazekas, M. Fleury, and M. Heisinger,
\emph{CaDiCaL, Kissat, Paracooba, Plingeling and Treengeling entering the SAT Competition 2020},
with current CaDiCaL documentation and sources.
\url{https://github.com/arminbiere/cadical}.

\bibitem{DRAT2014}
N. Wetzler, M. J. H. Heule, and W. A. Hunt Jr.,
\emph{DRAT-trim: efficient checking and trimming using expressive clausal proofs},
SAT 2014, LNCS 8561, 422--429.
\href{https://doi.org/10.1007/978-3-319-09284-3_31}{doi:10.1007/978-3-319-09284-3\_31}.

\end{thebibliography}
}

\end{document}
